\chapter{Introduction}
\label{ch:introduction}

This chapter establishes the research framework for SENTINEL, a two-tier cognitive architecture that automates threat detection and incident response in modern Security Operations Centers (SOCs). By analysing alert fatigue, deterministic rule blindness, and the computational constraints of Large Language Models (LLMs), it details the design and validation of a secure hybrid detection pipeline, delineating the objectives, research questions, scope, and quantitative methodology.

\section{Background and Motivation}
Modern enterprise networks have become cloud-native and porous, vastly expanding the attack surface. Adversaries span commodity malware to sophisticated multi-phase campaigns~\cite{mlids2016}, yet all converge on one chokepoint: the sheer volume of telemetry defenders must sift through. SOCs aggregate logs from NGFWs, IDS/IPS, and EDR into centralised SIEM platforms~\cite{siem2021} for normalisation and correlation. Such front-line filtering (line-rate IDS/IPS, edge ingestion) is deterministic and rigid: it curtails volume but can neither reason about an ambiguous alert nor explain a verdict, so the velocity of telemetry still produces a ``Big Data Paradox''~\cite{bigdata2019} that pushes contextual judgment back onto human analysts. This manifests as a severe \emph{alert-fatigue} crisis: analysts are inundated with low-fidelity, mostly false-positive notifications~\cite{alertfatigue2022}, dramatically increasing the risk of overlooking genuine intrusions.

Two further gaps compound this. First, signature-anchored filters~\cite{mlids2016} are structurally blind to novel variants (zero-day exploits, polymorphic malware) that bypass static rules with minor payload changes, exposing the need for a signature-independent statistical-outlier signal. Second, cost-driven log sampling severs the temporal linkages needed to reconstruct low-and-slow chains~\cite{apt_stealth2021}. Conventional SOAR does not close the reasoning gap either---it executes rigid pre-authored playbooks that cannot reason about or explain unscripted alerts~\cite{soar2025}.

LLMs promise exactly this missing contextual reasoning, but their direct use in live SOCs faces an integration paradox. Their inference latency is prohibitive for high-velocity logs due to the quadratic cost of self-attention~\cite{vaswani2017}; their natural-language interface exposes them to log-substrate prompt injection, where payloads embedded in fields such as User-Agent hijack the model's reasoning unless strict isolation is enforced~\cite{promptinjection2023}; and in regulated, air-gapped, or Zero-Trust settings~\cite{nist80207} telemetry cannot leave for a cloud model, so the reasoning tier must run locally. To resolve these, SENTINEL proposes a dual-tier model: a stateless Tier-1 filter at wire-speed, a learned gateway that sheds the confident majority of escalations, and a secure, locally hosted Tier-2 Agentic AI tier performing cognitively isolated reasoning. As a by-product, the stateful design also preserves low-and-slow chains and flags signature-less outliers---treated as foundational capabilities, not the central result.

\section{Research Objectives}
The general objective is to design, implement, and quantitatively evaluate the SENTINEL two-tier cognitive architecture combining low-latency statistical filtration with secure, context-aware agentic reasoning. Three specific objectives follow: (1)~engineer a wire-speed filter implementing Welford's online algorithm for $O(1)$ rolling variance and Z-scores, discarding benign traffic at the ingestion layer; (2)~design a stateful cognitive agent on LangGraph and a local quantised LLM, integrated with input--output guardrails that neutralise the adversarial vectors targeting LLM-integrated systems; and (3)~construct a dual retrieval-augmented grounding subsystem fusing MITRE ATT\&CK with NIST SP~800-61r2 to ground the agent's decisions and minimise hallucination.

\section{Research Questions}
Three questions guide the empirical validation. \textbf{RQ1:} how can a real-time filtration mechanism be designed on online statistical baselines to run at wire-speed and mitigate alert fatigue while preserving the temporal event chains of APT campaigns? \textbf{RQ2:} by what architectural methodology can local LLMs be integrated into an automated response pipeline, with secure input--output encapsulation, to minimise latency while mitigating adversarial attack vectors? \textbf{RQ3:} what empirical efficacy is gained by augmenting the agent's memory with an external dual knowledge-retrieval system based on Reciprocal Rank Fusion, as measured by classification accuracy and reduced hallucination?

\section{Scope and Object of Study}
\textbf{Object of study:} the two-tier cognitive architecture itself---the mechanism by which a stateless, deterministic Tier-1 filter at wire-speed is composed with a stateful, locally hosted Tier-2 agentic tier to automate SOC detection and response. This entails its constituents: high-velocity flow telemetry at OSI Layer~3--4 (the NetFlow five-tuple and CICFlowMeter flow features, with PCAP metadata rather than deep payload inspection); stateful agent orchestration (LangGraph); local quantised LLMs for air-gapped inference; and the adversarial vectors of log-substrate injection, delimiter smuggling, and knowledge-base poisoning. \textbf{Scope:} evaluation is restricted to CSE-CIC-IDS2018~\cite{ids2018} for general attacks and DAPT2020~\cite{dapt2020} for multi-stage APT, with zero-days modelled by injecting extreme Welford deviations into benign flows. The model is a local \textit{Gemma-2-9B-IT} quantised to Q6\_K via \texttt{llama.cpp} on a single consumer GPU, ensuring data sovereignty. Autonomous actions are limited to firewall-policy generation, reporting, and forensics, excluding automated physical disconnection.

\section{Research Methodology}
The thesis follows a Design Science Research approach under deductive, quantitative reasoning: it produces knowledge by building a working artifact (SENTINEL) and evaluating it against a real operational problem, settling every claim by measurement. It reasons statistically (separating normal traffic from anomalies), logically (structuring the agent's decision process), and adversarially (testing manipulation resistance). Concretely, the system is implemented as a two-tier prototype, exercised through end-to-end experiments that replay stored datasets and live traffic under adversarial inputs, and assessed on a five-dimensional framework (accuracy, performance, security, explainability, integrity). Controlled ablation isolates each component's contribution, and non-parametric tests confirm the differences are significant. The study draws on quantitative data (numerical measurements and outputs such as accuracy and latency) and qualitative data (the model's natural-language explanations), the latter scored via an automated framework~\cite{ragas2023} and an independent cross-family LLM judge so that reasoning quality is assessed on a comparable basis.

\section{Contributions of the Thesis}
The contributions are three-fold, each stated at the scope its evidence supports. \textbf{Architecturally}, the work validates a two-tier pipeline that makes a locally hosted, multi-second LLM operationally feasible in a SOC. The contribution is not a pre-filter in isolation but the \emph{composition}: a stateless $O(1)$ statistical gate and a learned load-shedding gateway absorb benign and confidently classifiable traffic at wire-speed and protect the network deterministically, while the costly cognitive tier is invoked asynchronously and only on the ambiguous minority. \textbf{In AI security}, it contributes a deterministic, structural defence against log-substrate prompt injection---Delimited Data Encapsulation, a Tier-1/Tier-2 consensus guard, and an HMAC-chained audit trail---whose core guarantee holds by construction rather than by training: a fixed-rule consensus guard cannot be argued down, whether or not a payload was seen before. \textbf{In orchestration}, it replaces rigid Runbook Automation with a stateful agent whose flow is an acyclic, tool-free state-graph on a single local quantised model, yielding a tier that is adaptive yet reproducible and forensically auditable under data-sovereignty constraints. Two further capabilities the architecture enables---signature-less anomaly escalation and emergent multi-day correlation---are offered as honestly bounded proofs of concept, with production maturation deferred to future work.

\section{Structure of the Thesis}
Chapter~\ref{ch:theoretical_background} establishes the theoretical foundations (LLM quantization, state-graph models, RAG, online statistics) and reviews related agentic-security work. Chapter~\ref{ch:system_design_implementation} details the system design: the two-tier filter, ML gateway, stateful agentic FSM, hybrid retrieval, guardrails, dashboard, and deployment. Chapter~\ref{ch:experiments_evaluation} presents the experiments, ablation, and statistical validation. Chapter~\ref{ch:conclusion} concludes with findings, limitations, and future directions.
